\documentclass{article}
\usepackage[a4paper, total={6.5in, 10in}]{geometry}
\setlength{\parindent}{0pt}
\usepackage{tcolorbox}
\tcbset{colback=white!94!black, colframe=black!60!white, coltitle = white, boxrule=0.8pt, arc=2mm}
\newtcolorbox{boxs}[2][]{title= #2,#1}
\usepackage{datetime2}
\usepackage[british]{babel}
\usepackage{amsfonts}
\usepackage{amsmath}
\usepackage{multicol}

\title{\textbf{Exam Prep Semester 1}}
\author{Robert O'Connell}
\date{\today}
\begin{document}
\maketitle
\section*{Single Variable Calculus}
\subsection*{Exam Layout}
3 hours long, 8 questions

Questions are similar to tutorial and homework problems

\subsection*{Content}
For a list of all content see Content.pdf

\subsection*{Tutorial Questions}
\begin{multicols}{2}
\begin{enumerate}
    \item \begin{itemize}
            \item Calculating domain and range
            \item \textbf{Showing bijectivity}
            \item Using rational root theorem and polynomial division
            \item Trigonometric addition formulae
        \end{itemize}
    \item \begin{itemize}
            \item Calculating inverse functions
            \item Simplfying trig. expression with inverses in them
            \item \textbf{Computing limits using epsilon-delta definition}
            \item Computing limits
            \item Determining when limits exist
        \end{itemize}
    \item \begin{itemize}
            \item Proving roots exist using Bolzano's theorem
            \item Determining when functions are continuous
            \item Computing limits at infinity/ when denominator is zero
            \item Using definition of derivative
            \item \textbf{Proving positvity at nearby points}
        \end{itemize}
    \item \begin{itemize}
            \item Computing derivatives by the rules
            \item Logarithmic diffrentiation
            \item Implicit diffrentiation
        \end{itemize}
    \item \begin{itemize}
            \item Using Rolle's theorem with Bolzano's
            \item \textbf{Using the mean value theorem to approximate values}
            \item Using L'hopital's rule to compute indeterminate limits
            \item Find values for which a function is increasing/ decreasing
            \item Find values for which a function is concave up/ down
        \end{itemize}
    \item \begin{itemize}
            \item Computing local minima/ maxima
            \item Computing global minima/ maxima on a closed interval
            \item Finding linear approximations to a function at a point
            \item Related rates of change
            \item Optimisation problems
            \item Estimating roots using Newton's method
        \end{itemize}
    \item \begin{itemize}
            \item Compute area under a curve
            \item Compute volume of a curve rotated
            \item Compute arc lengths of a graph
            \item Show a function is integrable using Riemann Sums
            \item Compute integrals using integration by parts and substitution
        \end{itemize}
    \item \begin{itemize}
            \item Computing integrals of powers of sine and cosine, tangent and cosecant
            \item Compute volume of a curve rotated
            \item Integration by parts
            \item Integration by substitution
            \item Reduction formulae
            \item Integration by trigonometric substitution
            \item Integration using partial fraction decomposition
        \end{itemize}
    \item \begin{itemize}
            \item Tests for convergence
            \item Geometric and p series
            \item Radius of convergence
            \item Differentiating power series
        \end{itemize}
\end{enumerate}
\end{multicols}

\newpage
\section*{Introduction to Programming}
\subsection*{Exam Layout}
2 hours long

\subsection*{Covered in Class}

\begin{multicols}{2}
    
\textbf{Number Systems}
\begin{itemize}
    \item Binary representations
    \item Converting integers to binary
    \item Converting fractions to binary
    \item Octal representations
    \item Hexadecimal representations
    \item Binary to octal/ hexadecimal
    \item Binary to decimal (Horners Method)
\end{itemize}

\textbf{Computer Representations of Numbers}
\begin{itemize}
    \item Fixed-point
    \item Floating-point
    \item Conversions to floating-point numbers
    \item Disadvantages of floating-point numbers
\end{itemize}

\textbf{Basic Programming}
\begin{itemize}
    \item Basics
    \item Loops
    \item Conditionals
    \item Functions
    \item Recursion
    \item Pseudocode
    \item Structures
\end{itemize}

\textbf{Logic}
\begin{itemize}
    \item AND, OR, NOT
    \item Boolean expressions
    \item Truth tables
\end{itemize}

\textbf{Proof Writing}
\begin{itemize}
    \item Structures of proofs
    \item Methods of proofs
\end{itemize}

\textbf{Computational Work}
\begin{itemize}
    \item Operation counts
    \item o-notations
\end{itemize}

\textbf{Covered Algorithms}
\begin{itemize}
    \item Estimate square roots (cube)
    \item Pi approximations
    \item Binary search (n-ary)
    \item Bisection root-finding
    \item Netwon's method
    \item Merge sort
\end{itemize}

\textbf{Linear Algebra}
\begin{itemize}
    \item Vector-Vector multiplication
    \item Matrix-Vector multiplication
    \item Matrix-Matrix multiplication
    \item Matrix structures
    \item Column/ row major
    \item Sparse matrices
    \item Strasen's method
\end{itemize}
\end{multicols}

\newpage
\section*{Introduction to Probability}
\subsection*{Exam Layout}
2 hours long, six questions

Questions will be quite similar to tutorial Questions

Just look at them :)



\end{document}